\begin{proposition}[The kernel on an open family of twelve-tooth smoothings]\label{inv:E2.5}
\StatementRecord{inv:E2.5}{proved}{}{the original marked comb construction of \cite{Viaclovsky2026Combs}; \cref{inv:E2.2}; \cref{cdp:st:E15}; the actual cup map \cref{invII:eq:E.6}}
Use the original marked section of the comb construction on
$\mathcal B=\mathbb P(3,4)$, with coarse base $B=\mathbb P^1$, the
order-three and order-four points $p_1,p_2$, and the hexagonal cusp $p_0$.
Its bundles are $E_0=G_0=\mathcal O\oplus\mathcal O(-1)$ and
$\mathcal L=\mathcal O(-5)$. There is a nonempty open family of
admissible positive-comb manifolds $X_{12}$ obtained by embedded
smoothings of twelve deformed twistor conics attached at distinct
ordinary points, fixing the three marked period planes, for which the
actual kernel in \eqref{invII:eq:E.6} satisfies
\[
 K_{12}\simeq\mathcal O_B(-8)\oplus\mathcal O_B(-9)^{\oplus2},
 \qquad \operatorname{coker}\kappa_1=0.
\]
On this family,
\[
 h^0(K_{12})=0,\qquad h^1(K_{12})=23,\qquad
 h^{1,1}(X_{12})=49,\qquad h^{1,2}(X_{12})=48.
\]
\end{proposition}

\begin{proof}
\noindent\textbf{Step 1. The actual seed kernel.}
Put $F_0=E_0^*$ and $D_0=\det G_0$. Write $x,y$ for the homogeneous
coordinates of weights $3,4$, respectively. Thus $y$ vanishes at the
order-three point. In row-$G$ and column-$F$ convention the actual
logarithmic period derivative has the form
\[
 T_0=\begin{pmatrix}0&ay\\by&cx\end{pmatrix},\qquad ab\ne0.
\]
Indeed its entries have weights $5,4,4,3$; the corresponding section
spaces are $0,\mathbb Cy,\mathbb Cy,\mathbb Cx$. Its invertible
logarithmic residue at the cusp forces $a,b$ to be nonzero.
For a tensor matrix $\eta$ in row-$F$ and column-$G$ convention define
\[
 \langle T,\eta\rangle
   =\sum_{i=1}^2(T_{2i}\eta_{i1}-T_{1i}\eta_{i2})
   =\operatorname{tr}(T\eta J),\qquad
 J=\begin{pmatrix}0&1\\-1&0\end{pmatrix}.
\]
This is the actual determinant-valued cup pairing. The seed row is
$(by,0,cx,-ay)$, with source weights $0,-1,1,0$ and target weight $4$.
Exact finite invariants give
\[
 \rho_*\mathcal O(4)=\mathcal O_B,
 \qquad \rho_*\mathcal O(1)=\rho_*\mathcal O(-1)=\mathcal O_B(-1).
\]
For the last two identities the invariant sheaves are lines with zero
$H^0$ and $H^1$, as follows from weighted homogeneous sections and
Serre duality with $K_{\mathcal B}=\mathcal O(-7)$.
The invariant section $y$ is the global coarse frame of the target.
Let $\sigma_x:\rho_*\mathcal O(1)\to\rho_*\mathcal O(4)$ be the
actual multiplication-by-$x$ map. The coarse row is
$(b,0,c\sigma_x,-a)$, which is onto. Its kernel isomorphism is
\[
 (u,v,w)\longmapsto
 (au/b-c\sigma_xv/b,\ w,\ v,\ u),
 \qquad \eta\longmapsto(\eta_{22},\eta_{21},\eta_{12}).
\]
Consequently $K_0=\mathcal O_B\oplus\mathcal O_B(-1)^2$.
This computation includes the possible $cx$ column. The determinant
$-ab y^2$ also proves invertibility of the actual handle derivative at
every ordinary attachment point.

\noindent\textbf{Step 2. The root family and the flat relative kernel.}
Choose twelve distinct ordinary points. The conics chosen below lie
in the fixed-marking embedded smoothing construction of
\cite[Proposition labelled \texttt{comb}]{Viaclovsky2026Combs}.
Here is the actual family lift at the finite points. In equivariant
frames $f_i,g_j$, the invariant generators are $w^{a_i^F}f_i$ and
$w^{a_j^G}g_j$. A coarse graph matrix $C(b,s)$ lifts by
\[
 Z_{ji}(w,s)=w^{a_j^G-a_i^F}C_{ji}(w^n,s).
\]
This follows by dividing its invariant graph vector by
$w^{a_i^F}$; multiplying back gives the inverse graph map.
At $p_1$ the order is $3$ and the exponent matrix is
$\left(\begin{smallmatrix}0&-1\\2&1\end{smallmatrix}\right)$.
At $p_2$ the order is $4$ and it is
$\left(\begin{smallmatrix}0&-3\\1&-2\end{smallmatrix}\right)$.
Since the entire period plane stays fixed along each marked section,
the required entries have the joint holomorphic divisions
$C_{12}=bD_{12}$, and also $C_{22}=bD_{22}$ at $p_2$.
Their lifts are respectively
\[
 w^2D_{12}(w^3,s),\qquad
 wD_{12}(w^4,s),\qquad w^2D_{22}(w^4,s).
\]
These are precisely the finite-character conditions of
\ref{cdp:st:E15}; all other exponents are nonnegative.
They construct the actual family Hodge and quotient bundles, with
their unchanged finite characters. At the cusp the chart in that
proposition is $(\log q)I/(2\pi i)+H(q,s)$ with $H$ jointly
holomorphic and bounded on a smaller parameter polydisc. Its residue
is unchanged. The same finite-cover actions, toric unit translations
and annular attachments give the compact family as in
\ref{cdp:st:E15}. In particular its coarse cusp cup row stays onto.

Choose a smoothing direction nonzero in every node-smoothing space.
Its regular total surface has coordinates $xy=s$ at each node,
with handle $y=0$ and tooth $x=0$. In an actual flat cohomology
frame write the period plane as the graph of $Z(x,y)$.
The relative dualizing generator $dx/x=-dy/y$ gives the derivative
\[
 \Theta=xZ_x-yZ_y.
\]
The relative cup takes values in
$D\otimes\omega_{\mathcal C/\Delta}(p_0)$, with the finite roots
along the marked sections understood. Put $H=Z_x(0,0)$ and
$B=Z_y(0,0)$. Our conic choices make the two functionals
$\langle H,-\rangle$ and $\langle B,-\rangle$ independent.
The cup-row entries then generate $(x,y)$ by their independent
linear terms and Nakayama's lemma. Invertible local operations
reduce the row to $(x,y,0,0)$, with free kernel generated by
$(-y,x,0,0)$ and the last two coordinate vectors.
Thus the relative kernel is locally free at every node.
Away from the nodes the central coarse row is onto; on the teeth
this follows from the conic calculation in Step 4. Its unit
coefficients persist after shrinking the family. Exact finite
invariants commute with parameter restriction in the character
frames just constructed. Finally $(x,y)$ is the unit ideal on a
noncentral fibre. We obtain a locally free coarse relative kernel
whose smooth restriction is the actual $K_{12}$ and whose smooth
coarse cokernel is zero.

\noindent\textbf{Step 3. The complete node condition.}
Let $\eta_h(x),\eta_t(y)$ be the two branch sections, with common
ordinary value $\eta_0$. Their value equations are
\[
 \langle H,\eta_0\rangle=\langle B,\eta_0\rangle=0.
\]
Write $Z=Z_h(x)+Z_t(y)-Z(0,0)+xyR(x,y)$. The coefficient of $xy$
in $xZ_x-yZ_y$ is zero, including the arbitrary mixed graph term.
The mixed coefficient in the kernel equation is therefore
\[
 \langle H,\eta_t'(0)\rangle
       -\langle B,\eta_h'(0)\rangle=0.                 \tag{E.12.1}\label{invII:eq:E12-node}
\]
This is also sufficient to all analytic orders. Over
$R_0=\mathbb C\{x,y\}$ let $\varphi:R_0^4\to R_0$ be the actual
row, with image $(x,y)$. A lift of the branch pair has
$\varphi(\widetilde\eta)=xyh$, and $h(0)$ is exactly the left side
of \eqref{invII:eq:E12-node}. Adding $xy\zeta$ changes $h$ by
$\varphi(\zeta)$, whose values range over $(x,y)$. The displayed
vanishing therefore gives an actual analytic correction with
$\varphi(\zeta)=-h$. The corrected lift is a section of the
total kernel with the prescribed restrictions. This proves that
the one mixed-jet equation is the entire remaining condition.
The kernel of the central row exceeds the central fibre of the
flat kernel by just one node copy of $\mathbb C$.
Changing a determinant or dualizing generator adds multiples of
the two zero value equations. Tensor-frame changes transport both
derivatives with their actual frame derivatives retained.

\noindent\textbf{Step 4. Conics and their actual two-jet map.}
In the flat decomposition $F_b\oplus G_b$ the through-point conics
under consideration have homogeneous inclusion and quotient
\[
 S_{[u:v]}=\binom{uI+vA}{vB},\qquad
 Q_{[u:v]}=(-vB,\ uI+vBAB^{-1}),\qquad \det B\ne0.
\]
Their product is zero. At $v\ne0$ the indicated bottom and left
blocks have full rank, and at $v=0$ the other blocks have full rank.
Thus they give the exact sequence with bundles
$F|_{C_i}=\mathcal O(-1)^2$ and $G|_{C_i}=\mathcal O(1)^2$.
The bottom-block chart recovers $u/v$ from the plane; moreover
$QS'=B$ near $y=v/u=0$. They are embedded immersed conics.
Set $P=I+yA$ and $C=BPB^{-1}$. Then
\[
 Z_t=yBP^{-1},\qquad S=(I;Z_t)P,
 \qquad Q=C(-Z_t,I),\qquad C Z_t'P=B.
\]
The tensor transition from the polynomial conic frames to graph
frames is $\eta_{\rm graph}=P\eta_{\rm conic}C^{-t}$.
The determinant factor is fixed by the exact identity
\[
 \langle CTP,P^{-1}\eta C^t\rangle
       =\det(C)\langle T,\eta\rangle,
\]
which follows from $C^tJC=(\det C)J$.
The ordinary tooth cup is the constant row $\langle B,-\rangle$
on $F\otimes G=\mathcal O^4$, so its kernel is $\mathcal O^3$.
Its global sections are constant in these frames. The derivative
forced by their value at the node is consequently
\[
 \eta_t'(0)=A\eta_0-\eta_0(BAB^{-1})^t.
\]
Substituting this exact frame derivative in \eqref{invII:eq:E12-node}
gives the two conditions on an ordinary handle-kernel section:
\[
 \begin{split}
 \langle B,\eta_h(0)\rangle&=0,\\
 \langle B,\eta_h'(0)\rangle
  &=\langle H,A\eta_h(0)-\eta_h(0)(BAB^{-1})^t\rangle.
 \end{split}                                        \tag{E.12.2}\label{invII:eq:E12-jets}
\]
The eight parameters $(A,B)$, modulo
$(A,B)\mapsto(rA+tI,rB)$ with $r\ne0$, are the actual
six-dimensional through-point conic chart: a polynomial basis
of its tautological $\mathcal O(-1)^2$ has constant blocks $I,0$
and linear blocks $A,B$. The two-jet kernel in
\eqref{invII:eq:E12-jets} is unchanged by these reparameterizations.
The opened conic lemma gives the same local deformation space and
a nonempty complex-open set of admissible compact conics.

\noindent\textbf{Step 5. Dominance onto the cyclic two-jet quotients.}
The perfect pairing $\langle-,-\rangle$ identifies
$B\mapsto\langle B,-\rangle|_{\ker\langle H,-\rangle}$ with the
quotient of a four-dimensional dual space by the line spanned by
$H$. It gives every first covector, subject to a dense open restriction.
For the second coefficient the following frame map is explicit.
Under $F'=FU$, $G'=GV$ one has
\[
 T'=V^{-1}TU,\quad \eta'=U^{-1}\eta V^{-t},\quad
 \langle T',\eta'\rangle=(\det V)^{-1}\langle T,\eta\rangle,
 \quad A'=U^{-1}AU,\quad B'=V^{-1}BU.
\]
Choose an eigenvector matrix $R$ for $H^{-1}B$ and put $U=R,V=HR$.
The actual derivative becomes $I$, the actual conic direction becomes
$\operatorname{diag}(\beta_1,\beta_2)$, and the tensor becomes
$R^{-1}\eta H^{-t}R^{-t}$. Work on the nonempty open set
\[
 \beta_1\beta_2(\beta_1-\beta_2)(\beta_1+\beta_2)\ne0.
\]
The two value equations then give
$\eta_0=\operatorname{diag}(u,v)$. For
$A=\left(\begin{smallmatrix}p&q\\r&s\end{smallmatrix}\right)$,
the second coefficient in these actual frames is
\[
 \langle I,A\eta_0-\eta_0(BAB^{-1})^t\rangle
   =r(1+\beta_2/\beta_1)u-q(1+\beta_1/\beta_2)v.
\]
Its Jacobian in $r,q$ has determinant
$-(\beta_1+\beta_2)^2/(\beta_1\beta_2)$, so the two coefficients
are arbitrary. The scalar part of $A$ contributes zero, as required
by the source reparameterization.
In any fixed local handle frame the resulting quotient is
\[
 \eta\longmapsto\ell(\eta(0))+
 t\bigl(\ell(\eta'(0))+\psi(\eta(0))\bigr)\pmod{t^2},
 \qquad \ell\ne0,
\]
where $\psi$ is arbitrary modulo $\ell$. The actual derivative of
the chosen handle frame translates $\psi$. This proves dominance
from actual conics to the four-dimensional cyclic length-two
quotient space, with its complete local frame map.

\noindent\textbf{Step 6. An attained balanced open set.}
Let $K'$ be the subsheaf of $K_0$ defined by the twelve quotients
\eqref{invII:eq:E12-jets}. Their parameter space is irreducible.
There is an exact witness for its balanced locus: assign four of
the distinct nodes to each summand of
$K_0=\mathcal O\oplus\mathcal O(-1)^2$, and take the coordinate
quotient modulo the square of the local parameter at each assigned
point. Multiplication by each product of four squared linear
factors gives its actual kernel inclusion. The resulting bundle is
$\mathcal O(-8)\oplus\mathcal O(-9)^2$, so
\[
 H^0(K'(7))=0,\qquad H^1(K'(8))=0.                    \tag{E.12.3}\label{invII:eq:E12-open}
\]
These conditions define a nonempty Zariski-open subset by
semicontinuity. Step 5 pulls it back to a nonempty open subset of
actual conic data; no particular coordinate witness is asserted to
be an admissible conic.

The original handle normal bundle has degrees $(-1,0,0,1)$.
With the three fixed marks, positive modifications distributed as
$(4,3,3,2)$ give a globally generated modified normal bundle with
zero $H^1$. The complete normal-sheaf proof in
\cite[Proposition \texttt{comb}, equation \texttt{originalnormal}
and the subsequent twelve-tooth bound]{Viaclovsky2026Combs}
shows that this is a second nonempty Zariski-open condition in the
actual tangent directions. The conic tangent map is dominant, so
both conditions pull back to dense open subsets of the same
irreducible conic chart. Their intersection meets the nonempty
complex-open set of admissible compact conics.
For these choices the cited fixed-marking deformation theorem
gives an embedded smoothing with a direction nonzero at every
node. The curve has genus zero and degree one over $B$, so each
smooth curve is an actual section. It fixes all three marked
period planes and has the root and compact filling of Step 2.

\noindent\textbf{Step 7. Collapse, cohomology and the smooth bundle.}
Let $K_C$ denote the central fibre of the locally free relative
kernel. Its handle restriction maps into $K_0$ by the first
elementary modification at each node. On a tooth its image in the
ordinary $\mathcal O^3$ kernel is cut out by the additional nonzero
value functional $\langle H,-\rangle$, and hence
\[
 K_C|_{C_i}=\mathcal O(-1)\oplus\mathcal O^2.
\]
This follows by choosing that functional as the first coordinate:
the first coordinate vanishes at the node and the other two are
constant. Its $H^1$ vanishes, and its two global sections evaluate
injectively into the three-dimensional relative-kernel node fibre.
Their image retains exactly the derivative constraint of Step 3.

Collapse the teeth by the actual projection $\pi:C\to B$.
The normalization sequence for $K_C$, the zero tooth $H^1$ and the
surjective local evaluation of the handle onto each whole node
fibre give $R^1\pi_*K_C=0$. A handle section extends uniquely over
each tooth precisely by the two equations in
\eqref{invII:eq:E12-jets}, including the sufficient mixed condition.
Thus restriction gives the exact sheaf identification
$\pi_*K_C=K'$. For the actual pullback line from the base this gives
\[
 H^j(C,K_C\otimes\pi^*\mathcal O_B(k))
       =H^j(B,K'(k)),\qquad j=0,1.
\]
The projection of the embedded family supplies this single line
throughout the smoothing: it is trivial on the central teeth and
is the required $\mathcal O_B(k)$ on each smooth section.
Flatness and semicontinuity therefore transport
\eqref{invII:eq:E12-open} to all sufficiently small noncentral
members of the chosen smoothing.
The central Euler characteristic is
\[
 (1-12)+12\cdot2-12\cdot3=-23,
\]
so its rank-three degree is $-26$. Equivalently, the smooth
surjective cup row has source degree $-2$ and target degree $24$.

For completeness, these two vanishings determine the bundle by
explicit maps. Riemann--Roch gives $h^0(K_{12}(8))=1$. Its section
is nowhere zero, since a zero would give a section of $K_{12}(7)$.
The quotient of $K_{12}(8)$ by this copy of $\mathcal O$ has rank
two, degree $-2$ and zero $H^0,H^1$. Evaluation of its twist by
one is onto at every point; Riemann--Roch gives two source
dimensions, so that evaluation is an isomorphism. The quotient
is $\mathcal O(-1)^2$, and the extension by $\mathcal O$ splits
because $H^1(\mathcal O(1))^2=0$. Twisting back proves the stated
kernel isomorphism. The two vanishings and surjectivity are open
in the actual nearby fixed-marking section family, giving the
nonempty open family in the statement. Formula
\eqref{invII:eq:E.5} supplies its two Hodge numbers.
\end{proof}

\paragraph{Relation to the original construction.}
The conic geometry, fixed-marking embedded smoothing and sufficient
twelve-tooth bound are those of
\cite[Lemma \texttt{twistor}, Proposition \texttt{comb}, and equation
\texttt{originalnormal}]{Viaclovsky2026Combs}.
The proof above computes the actual flat cup kernel and its two-jet
modifications, then transports their cohomology through that
smoothing. It determines the stated open twelve-tooth family in the
remaining-Hodge-number question of \cite[Section~10(2)]{Viaclovsky2026Combs}.
The general and special-member splitting problem remains in
\ref{inv:E2.3}.
